% apple-pay-passkit.tex — Apple Pay in apps: merchant ID + Payment Processing Certificate, entitlement,
% PKPaymentRequest, the authorization-controller delegate flow, the encrypted token (DPAN + cryptogram),
% PSP decryption, merchant tokens for recurring / deferred / auto-reload, Wallet order tracking,
% Apple Pay vs In-App Purchase, Apple Pay on the web.
% IAP itself is the storekit2 folio; Wallet passes (.pkpass) are the wallet-passes folio.
% Sources (checked 2026-10-05):
%   developer.apple.com/documentation/passkit (PKPaymentRequest, PKMerchantCapability,
%   PKRecurringPaymentRequest, PKPaymentOrderDetails, PayWithApplePayButton,
%   PKPaymentAuthorizationControllerDelegate)
%   developer.apple.com/documentation/passkit/payment-token-format-reference (EC_v1/RSA_v1, 5 min, ECI)
%   developer.apple.com/app-store/review/guidelines (3.1.1, 3.1.3(e), 3.2.1(vi), 4.9)
%   WWDC24 "What's new in Wallet and Apple Pay" (Apple Pay in third-party browsers)
%   apple.com/newsroom (Apple Pay launch with iOS 8.1, 2014-10-20; Tap to Pay on iPhone, iOS 15.4)
% @labels: area=ios-platform kind=api platform=ios topic=platform-apis,security discussion=13
% @tags: apple-pay, passkit, pkpaymentrequest, payment-token, dpan, merchant-token, payment-processing-certificate, psp, three-d-secure, in-app-purchase-boundary, recurring-payments, wallet-orders
\documentclass[8pt]{extarticle}
\usepackage{printup-sheet}
\usepackage{array}

\lstdefinelanguage{SwiftSheet}{
  morekeywords={protocol,class,final,struct,enum,func,var,let,static,some,init,
    if,else,return,guard,self,nil,private,true,false,in,try,await,async,throws,
    switch,case,default,for,escaping,Void,String,Bool,Int},
  sensitive=true, morecomment=[l]{//}, morestring=[b]",
  literate={->}{{\hbox{-}\hbox{>}}}2 {==}{{\hbox{=}\hbox{=}}}2 {??}{{\hbox{?}\hbox{?}}}2}
\lstset{basicstyle=\ttfamily\scriptsize, aboveskip=2pt, belowskip=2pt}
\newcommand\ct[1]{\texttt{#1}}
\newcommand\hd[1]{\par\vspace{3pt}\noindent{\bfseries\color{sheetBlue}#1}\par\vspace{1pt}}

\tikzset{
  st/.style={box, font=\scriptsize, text width=27mm, minimum height=10mm, inner sep=2pt},
  dev/.style={st, draw=sheetBlue, fill=sheetBlue!8},
  apl/.style={st, draw=sheetGrey, fill=black!5},
  you/.style={st, draw=sheetGreen, fill=sheetGreen!10},
  psp/.style={st, draw=sheetOrange, fill=sheetOrange!8},
  lbl/.style={font=\tiny, text=black!75, inner sep=1pt, align=center},
}

\begin{document}

\sheettitle{Apple Pay in apps · PassKit payments}{ios-platform · folio}

\oneliner{Apple Pay pays for \textbf{physical goods and services} with a Wallet card: the app builds a
\ct{PKPaymentRequest}, the system sheet authenticates, the \textbf{Secure Element} emits a
\textbf{DPAN + one-time cryptogram}, Apple re-encrypts it to \textbf{your merchant key}, and the app
forwards the opaque \ct{token.paymentData} to your server $\to$ \textbf{PSP}. The app never sees a card
number. Digital content used in the app is \textbf{In-App Purchase}, not Apple Pay.}

\vspace{2pt}
\noindent\begin{tikzpicture}[sheet]
  % row 1: device + Apple, left to right
  \node[dev] (app) at (1.45,1.45) {\textbf{1 · App}\\\ct{PKPaymentRequest}\\+ \ct{present()}};
  \node[dev] (sh) at (4.75,1.45) {\textbf{2 · Payment sheet}\\(system process)\\Face ID / double-click};
  \node[dev, draw=sheetRed, fill=sheetRed!6] (se) at (8.05,1.45) {\textbf{3 · Secure Element}\\\textbf{DPAN} + one-time\\cryptogram (EMV/3DS)};
  \node[apl] (as) at (11.35,1.45) {\textbf{4 · Apple Pay servers}\\re-encrypt to the\\\textbf{merchant's} public key};
  \node[dev] (tok) at (14.65,1.45) {\textbf{5 · App} gets \ct{PKPayment}\\\ct{token.paymentData}\\(opaque JSON blob)};
  % row 2: right to left
  \node[you] (srv) at (14.65,0) {\textbf{6 · Your server}\\amount, order id,\\idempotency key};
  \node[psp] (pspn) at (11.35,0) {\textbf{7 · PSP} (Stripe, Adyen)\\decrypts with the\\\textbf{processing cert}'s key};
  \node[psp] (net) at (8.05,0) {\textbf{8 · Card network}\\token service maps\\\textbf{DPAN $\to$ FPAN}};
  \node[psp] (iss) at (4.75,0) {\textbf{9 · Issuer}\\checks cryptogram\\approve / decline};
  \node[dev] (done) at (1.45,0) {\textbf{10 · App} \ct{handler(.success)}\\sheet shows \checkmark\\\ct{DidFinish} $\to$ dismiss};
  \draw[flow] (app) -- (sh); \draw[flow] (sh) -- (se); \draw[hot] (se) -- (as); \draw[hot] (as) -- (tok);
  \draw[hot] (tok) -- node[lbl, right]{HTTPS} (srv);
  \draw[hot] (srv) -- (pspn); \draw[hot] (pspn) -- (net); \draw[hot] (net) -- (iss);
  \draw[flow, draw=sheetGreen] (iss) -- (done);
  % annotations
  \node[lbl, text=sheetRed, anchor=north] at (8.05,0.88) {the real card number (FPAN) is never on the phone or at the merchant};
  \node[lbl, text=sheetBrown, anchor=south] at (8.05,2.05) {delegate callbacks (shipping contact / method / coupon) update the totals before step 3};
  \node[lbl, text=sheetBlue, anchor=north] at (13.0,-0.62) {the app never decrypts: the private key lives at the PSP (or your HSM)};
  \node[lbl, text=sheetGrey, anchor=north] at (3.1,-0.62) {result travels back via PSP + server; decline $\to$ \ct{.failure} + errors, sheet stays up};
\end{tikzpicture}

\begin{multicols}{2}
\raggedright\setstretch{1.0}

\hd{Setup}
\begin{itemize}
  \item \textbf{Merchant ID} \ct{merchant.com.example.app} (portal). \textbf{Payment Processing
        Certificate} per merchant ID: the \textbf{PSP makes the CSR} and keeps the private key — the
        token is encrypted to it. The ID never expires; certificates do.
  \item Apple Pay capability $\to$ entitlement \ct{com.apple.developer.in-app-payments} (merchant IDs).
  \item \textbf{Web only}: \textbf{Merchant Identity Certificate} (TLS client cert) + a verified domain.
\end{itemize}

\hd{The request}
\begin{itemize}
  \item \ct{countryCode} = the \textbf{merchant's} country; \ct{currencyCode};
        \ct{supportedNetworks}; \ct{merchantCapabilities = .threeDSecure} \textbf{always}.
  \item \ct{paymentSummaryItems}: lines, then \textbf{last = grand total} labelled with your
        \textbf{business name} (``Pay EXAMPLE CO''). \ct{NSDecimalNumber}; \ct{.pending} if unknown yet.
  \item \ct{requiredShippingContactFields}, \ct{shippingMethods}, \ct{supportsCouponCode}.
        Before auth, \ct{didSelectShippingContact} gets a \textbf{redacted} address (city, postcode,
        country) — enough for tax; full contact only in the authorized \ct{PKPayment}.
\end{itemize}

\hd{Example — request, present, authorize}
\begin{lstlisting}[language=SwiftSheet]
let r = PKPaymentRequest()
r.merchantIdentifier = "merchant.com.example.app"
r.countryCode = "US"; r.currencyCode = "USD"
r.supportedNetworks = [.visa, .masterCard, .amex]
r.merchantCapabilities = .threeDSecure
r.paymentSummaryItems = [PKPaymentSummaryItem(label: "EXAMPLE CO",
    amount: NSDecimalNumber(string: "12.00"))]   // last = total
let c = PKPaymentAuthorizationController(paymentRequest: r)
c.delegate = self; c.present(completion: nil)
// delegate: charge server-side, report exactly once
func paymentAuthorizationController(_ c: PKPaymentAuthorizationController,
    didAuthorizePayment p: PKPayment,
    handler done: @escaping (PKPaymentAuthorizationResult) -> Void) {
  api.charge(p.token.paymentData) { ok in
    done(.init(status: ok ? .success : .failure, errors: nil)) } }
// required; also called on cancel
func paymentAuthorizationControllerDidFinish(_ c: ...) { c.dismiss() }
\end{lstlisting}

\hd{Button \& availability}
\ct{canMakePayments()} = device capable, not restricted; \ct{canMakePayments(usingNetworks:)} = a
matching card exists. No card $\to$ \ct{PKPaymentButton(type: .setUp)}. SwiftUI:
\ct{PayWithApplePayButton} (iOS 16). Apple's button only (4.9, HIG).

\hd{Apple Pay or In-App Purchase?}
{\footnotesize
\begin{tabular}{@{}>{\raggedright\arraybackslash}p{37mm}>{\raggedright\arraybackslash}p{36mm}@{}}
\toprule
digital content / subs used \emph{in} the app & \textbf{IAP} (3.1.1, StoreKit) \\
physical goods, services used \emph{outside} & \textbf{Apple Pay} / card — IAP not allowed (3.1.3(e)) \\
nonprofit donations & must offer Apple Pay (3.2.1(vi)) \\
accept cards on your iPhone & \textbf{Tap to Pay on iPhone} (ProximityReader) \\
\bottomrule
\end{tabular}}

\columnbreak

\hd{The token — what your server receives}
\ct{paymentMethod} (``Visa 1234'', network, debit/credit — UI only) · \ct{transactionIdentifier} ·
\ct{paymentData} = JSON \{\ct{version} (\ct{EC\_v1} ECC · \ct{RSA\_v1} RSA), \ct{data}, \ct{signature},
\ct{header} \{\ct{ephemeralPublicKey}, \ct{publicKeyHash}, \ct{transactionId}\}\}. \textbf{Decrypted} (PSP):
\ct{applicationPrimaryAccountNumber} = \textbf{DPAN}, expiry, amount, currency,
\ct{onlinePaymentCryptogram} + optional \ct{eciIndicator} — if present it \textbf{must} reach the
processor, else the transaction fails. CMS signing time vs transaction time $>$ \textbf{5 min} apart
$\to$ possible replay.

\hd{DPAN vs FPAN vs merchant token}
\textbf{FPAN} = printed card number. \textbf{DPAN} = network token in \emph{this device's} SE — gone
when the card is removed or the phone wiped. \textbf{Merchant token (MPAN)} = bound to you + the user's
Apple Account, survives a new phone. Get one via (iOS 16) \ct{recurringPaymentRequest}
(\ct{PKRecurringPaymentRequest}: description, \ct{regularBilling}, \ct{managementURL},
\ct{tokenNotificationURL}) · \ct{automaticReloadPaymentRequest} (top-ups) ·
\ct{deferredPaymentRequest} (16.4, hotels). \ct{multiTokenContexts} = one token per merchant
(marketplace); not combinable with recurring / auto-reload.

\hd{After payment \& on the web}
\ct{PKPaymentAuthorizationResult.orderDetails} = \ct{PKPaymentOrderDetails} (order type ID, order ID,
web service URL, auth token; iOS 16) $\to$ Wallet tracks the \textbf{order}. Web:
\ct{ApplePaySession} / Payment Request API; your server does \textbf{merchant validation} with the
Merchant Identity cert. iOS 18: non-Safari browsers show a \textbf{code the iPhone scans}.

\hd{Traps}
\begin{itemize}
  \trap{Digital upgrade via Apple Pay $\to$ rejected; a pizza via IAP $\to$ wrong too.}
  \trap{\ct{Double} for money; last item ``Total'' instead of the business name.}
  \trap{Decrypting \ct{paymentData} in the app = shipping the private key.}
  \trap{Storing a DPAN for subscriptions — request a merchant token.}
  \trap{Hiding the button when no card: show \ct{.setUp} instead.}
  \trap{Apple Pay (you \emph{charge}) $\neq$ Tap to Pay (you \emph{accept}) $\neq$ HCE.}
\end{itemize}

\hd{Timeline}
\begin{tikzpicture}[sheet]
  \tikzset{yr/.style={font=\bfseries\scriptsize, text=sheetBlue},
           ev/.style={font=\tiny, align=center, text=black!85, text width=13mm, anchor=north}}
  \draw[thick, sheetGrey, ->] (0,0) -- (8.3,0);
  \foreach \x/\y/\t in {
      0.6/iOS 8.1/{Apple Pay\\(US, 2014)},
      1.95/iOS 10/{Apple Pay\\on the web},
      3.3/iOS 15.4/{Tap to Pay\\on iPhone},
      4.65/iOS 16/{merchant tokens\\order tracking\\SwiftUI button},
      6.0/iOS 16.4/{deferred\\payments},
      7.35/iOS 18/{other browsers:\\scan a code}}{
    \fill[sheetBlue] (\x,0) circle (0.6mm);
    \node[yr, above=1pt] at (\x,0) {\y};
    \node[ev] at (\x,-0.12) {\t};
  }
\end{tikzpicture}

\hd{Remember}
\textbf{``ID, cert, sheet, SE, blob, PSP.''} Phone makes a one-time proof; only the PSP opens it; only
the network knows the card.

\end{multicols}

\noindent{\footnotesize\color{sheetGrey}\textit{Related:} storekit2 (IAP) · store-policy-and-deadlines (3.1.x) ·
wallet-passes · nfc-fundamentals (SE vs HCE) · keychain-secure-enclave · tls-pki (client certs)}

\end{document}
